\SourcePage{405}{408}
\section{Matrix elements for a diatomic molecule}

This section gives several general formulas for matrix elements of physical quantities pertaining to a diatomic molecule. We first consider matrix elements for transitions between zero-spin states.

Let $\mathbf A$ be a vector physical quantity that characterizes the molecule when its nuclei are fixed (for example, its electric or magnetic dipole moment). We first consider this quantity in the $\xi\eta\zeta$ coordinate system, which rotates with the molecule and whose $\zeta$ axis lies along the molecular axis. The molecule's angular momentum relative to this frame (that is, the electronic angular momentum $\mathbf L$) is not conserved in its entirety, whereas its $\zeta$ component is conserved. The selection rules associated with the quantum number $L_\zeta=\Lambda$ therefore continue to apply (they coincide with the selection rules for $M$ in \S~29). Hence the vector can have nonzero matrix elements only of the following kinds:

\begin{equation}
\langle n'\Lambda|A_\zeta|n\Lambda\rangle,\quad
\langle n'\Lambda|A_\xi+iA_\eta|n,\Lambda-1\rangle,\quad
\langle n',\Lambda-1|A_\xi-iA_\eta|n\Lambda\rangle
\tag{87.1}
\end{equation}
(here $n$ labels the electronic terms for a specified $\Lambda$).

When both terms are $\Sigma$ terms, the selection rule associated with reflection symmetry in a plane containing the molecular axis must also be taken into account. Under such a reflection, the $\zeta$ component of an ordinary (polar) vector remains unchanged, whereas that of an axial vector reverses sign. It follows that, for a polar vector, $A_\zeta$ has nonzero matrix elements only in the transitions $\Sigma^+\to\Sigma^+$ and $\Sigma^-\to\Sigma^-$, while for an axial vector this occurs for $\Sigma^+\to\Sigma^-$. We do not discuss the components $A_\xi$, $A_\eta$, since transitions without a change in $\Lambda$ are altogether impossible for them.

For a molecule made up of identical atoms, there is an additional selection rule involving parity.
The components of a polar vector reverse sign under inversion.
Its matrix elements can therefore be nonzero only for transitions between states of opposite parity (the reverse applies to an axial vector).
In particular, every diagonal matrix element of a polar-vector component vanishes identically.


The relation between the matrix elements (87.1) and the matrix elements of the same vector in the space-fixed coordinate system $xyz$ follows from the general formulas derived below (in \S~110) for an arbitrary axially symmetric physical system.

On separating the dependence on the quantum number $M_K$ that is common to any vector (the $z$ projection of the total angular momentum
\SourcePage{406}{409}
of the molecule, $\mathbf K$), one is left with the reduced matrix elements $\langle n'K'\Lambda'\|A\|nK\Lambda\rangle$.
Their relation to the matrix elements (87.1) is specified by formula (110.7), with $k=k'=1$ (as appropriate to a vector) and the corresponding relabeling of the quantum numbers (recall that, by (82.4), $\Lambda$ equals the $\zeta$ component of the total angular momentum $\mathbf K$).
Using relation (107.1) between the components of a rank-one spherical tensor and the Cartesian components of a vector, together with the values of the $3j$ symbols in Table~9 (p.~530), we obtain the following formulas for matrix elements diagonal in $\Lambda$:

\begin{equation}
\begin{split}
\langle n'K\Lambda\|A\|nK\Lambda\rangle&=\Lambda\sqrt{\frac{2K+1}{K(K+1)}}\langle n'\Lambda|A_\zeta|n\Lambda\rangle,\\
\langle n',K-1,\Lambda\|A\|nK\Lambda\rangle&=i\sqrt{\frac{K^2-\Lambda^2}{K}}\langle n'\Lambda|A_\zeta|n\Lambda\rangle
\end{split}
\tag{87.2}
\end{equation}
and, for elements not diagonal in $\Lambda$,
\begin{equation}
\begin{split}
\langle n'K\Lambda\|A\|nK,\Lambda-1\rangle
&=\left[\frac{(2K+1)(K+\Lambda)(K-\Lambda+1)}{4K(K+1)}\right]^{1/2}\\[-2pt]
&\qquad\times\langle n'\Lambda|A_\xi+iA_\eta|n,\Lambda-1\rangle,\\
\langle n'K\Lambda\|A\|n,K-1,\Lambda-1\rangle
&=i\left[\frac{(K+\Lambda)(K+\Lambda-1)}{4K}\right]^{1/2}\\[-2pt]
&\qquad\times\langle n'\Lambda|A_\xi+iA_\eta|n,\Lambda-1\rangle,\\
\langle n',K-1,\Lambda\|A\|nK,\Lambda-1\rangle
&=i\left[\frac{(K-\Lambda)(K-\Lambda+1)}{4K}\right]^{1/2}\\[-2pt]
&\qquad\times\langle n'\Lambda|A_\xi+iA_\eta|n,\Lambda-1\rangle.
\end{split}
\tag{87.3}
\end{equation}
The other nonzero elements follow from those displayed above upon using the Hermiticity relations for the reduced matrix elements,
\[
\langle nK\Lambda\|A\|n'K'\Lambda'\rangle=\langle n'K'\Lambda'\|A\|nK\Lambda\rangle^*,
\]
and for the matrix elements in the $\xi\eta\zeta$ system,
\[
\langle n\Lambda|A_\xi-iA_\eta|n'\Lambda'\rangle=\langle n'\Lambda'|A_\xi+iA_\eta|n\Lambda\rangle^*,\qquad
\langle n\Lambda|A_\zeta|n'\Lambda'\rangle=\langle n'\Lambda'|A_\zeta|n\Lambda\rangle^*.
\]

\SourcePage{407}{410}
We record separately the formulas for matrix elements of the vector $\mathbf A=\mathbf n$, the unit vector along the molecular axis. In this case simply $A_\xi=A_\eta=0$, $A_\zeta=1$, so in the $\xi\eta\zeta$ system only the diagonal elements are nonzero: $\langle n\Lambda|A_\zeta|n\Lambda\rangle=1$. The reduced matrix elements are diagonal in every index except $K$; displaying only that index, we have
\begin{equation}
\langle K\|n\|K\rangle=\Lambda\sqrt{\frac{2K+1}{K(K+1)}},\qquad
\langle K-1\|n\|K\rangle=i\sqrt{\frac{K^2-\Lambda^2}{K}}.
\tag{87.4}
\end{equation}
(H.~Hönl, F.~London, 1925). For $\Lambda=0$, these formulas give
\[
\langle K\|n\|K\rangle=0,\qquad \langle K-1\|n\|K\rangle=i\sqrt K,
\]
which is precisely what is expected for the matrix elements of a unit vector in motion in a centrally symmetric field (see (29.14)).

We now indicate how the formulas obtained above are to be modified for transitions between states of nonzero spin. Here it is essential whether the states belong to case $a$ or to case $b$.

If both states belong to case $a$, the changes in the formulas are essentially only notational. The quantum numbers $K$ and $M_K$ do not exist; in their place one has the total angular momentum $J$ and its $z$ projection $M_J$. In addition, the numbers $S$ and $\Omega=\Lambda+\Sigma$ are introduced, so that the reduced matrix elements are written as
\[
\langle n'J'S'\Omega'\Lambda'\|A\|nJS\Omega\Lambda\rangle.
\]
Let $\mathbf A$ be any orbital vector (that is, one independent of spin). Its operator commutes with the spin operator $\widehat{\mathbf S}$, and its matrix is therefore diagonal in the quantum numbers $S$ and $S_\zeta=\Sigma$; hence the quantum number $\Omega=\Lambda+\Sigma$ changes together with $\Lambda$ (that is, $\Omega'-\Omega=\Lambda'-\Lambda$). Equations (87.2)--(87.4) change only in that indices are added to the matrix elements and $K,\Lambda$ must be replaced by $J,\Omega$ in all the other factors. For example, the first of equations (87.2) is replaced by
\[
\langle n'JS\Omega\Lambda\|A\|nJS\Omega\Lambda\rangle
=\Omega\sqrt{\frac{2J+1}{J(J+1)}}\langle n'\Omega\Lambda|A_\zeta|n\Omega\Lambda\rangle
\]
(the diagonal index $S$ has been omitted).

Now let $\mathbf A=\mathbf S$. Since the spin operator commutes both with the orbital angular momentum and with the Hamiltonian, its matrix is diagonal in $n,\Lambda$. It is not, however, diagonal in $S$ and $\Sigma$ (or $\Omega$). For

\SourcePage{408}{411}
the transitions $S,\Sigma\to S',\Sigma'$, the matrix elements of the components $A_\xi,A_\eta,A_\zeta$ are determined by (27.13), with $S,\Sigma$ written in place of $L,M$. The transformation to the coordinate system $xyz$ is then performed using (87.2), (87.3), with $K,\Lambda$ replaced by $J,\Omega$. In this manner, for example, we obtain
\begin{equation}
\begin{split}
\langle J\Omega\|S\|J,\Omega-1\rangle
&=\left[\frac{(2J+1)(J+\Omega)(J-\Omega+1)}{4J(J+1)}\right]^{1/2}
\langle\Omega|S_\xi+iS_\eta|\Omega-1\rangle\\
&=\left[\frac{(2J+1)(J+\Omega)(J-\Omega+1)(S+\Sigma)(S-\Sigma+1)}{4J(J+1)}\right]^{1/2}.
\end{split}
\tag{87.5}
\end{equation}
(the diagonal indices $n,S,\Lambda$ have been omitted).

Next suppose that both states belong to case $b$, and that $\mathbf A$ is an orbital vector. The matrix elements are calculated in two stages. First we consider the rotating molecule without allowing for the coupling of $\mathbf S$ and $\mathbf K$; the matrix elements are diagonal in $S$ and are given by the same formulas (87.2), (87.3). At the second stage, $\mathbf K$ is coupled with $\mathbf S$ to form the total angular momentum $\mathbf J$, and the transformation to the new matrix elements is made by means of the general formulas (109.3) (with $K,S,J$ playing the parts of $j_1,j_2,J$ in those formulas). Thus, for elements diagonal in $J,K,\Lambda$, we first find
\[
\langle n'JK\Lambda\|A\|nJK\Lambda\rangle
=(-1)^{K+J+S+1}(2J+1)
\begin{Bmatrix}K&J&S\\J&K&1\end{Bmatrix}
\langle n'K\Lambda\|A\|nK\Lambda\rangle
\]
and then, using the value of the $6j$ symbol from Table~10 (p.~542) and the reduced matrix element from (87.2), we finally obtain
\[
\langle n'JK\Lambda\|A\|nJK\Lambda\rangle
=\Lambda\left[\frac{2J+1}{J(J+1)}\right]^{1/2}
\frac{J(J+1)+K(K+1)-S(S+1)}{2K(K+1)}
\langle n'\Lambda|A_\zeta|n\Lambda\rangle.
\]

Matrix elements for transitions between states belonging respectively to cases $a$ and $b$ are calculated in the same manner; we shall not pursue that calculation here.


\ProblemHeading[1] Determine the Stark splitting of the terms of a diatomic molecule having a permanent dipole moment, when the term belongs to case $a$.

\SourcePage{409}{412}
\SolutionHeading The energy of a dipole $\mathbf d$ in an electric field $\boldsymbol{\mathcal E}$ is $-\mathbf d\boldsymbol{\mathcal E}$. By symmetry, the dipole moment of a diatomic molecule is plainly directed along its axis: $\mathbf d=d\mathbf n$ (where $d$ is constant). Choosing the field direction as the $z$ axis, we obtain the perturbation operator in the form $-dn_z\mathcal E$.

Evaluating the diagonal matrix elements of $n_z$ from the formulas derived in the text, we find that in case $a$ the splitting of the levels is given by\footnote{This may appear to conflict with the general statement (see \S~76) that there is no linear Stark effect. In fact, there is of course no conflict, because here the linear effect is associated with the twofold degeneracy of the levels with $\Omega\ne0$; the formula obtained is consequently applicable provided that the energy of the Stark splitting is large compared with the energy of the so-called $\Lambda$-doubling (see \S~88).}

\[
\Delta E=-\mathcal E dM_J\frac{\Omega}{J(J+1)}.
\]

\ProblemHeading[2] The same problem for a term belonging to case $b$ (with $\Lambda\ne0$).

\SolutionHeading In the same manner we find
\[
\Delta E=-\mathcal E dM_J\Lambda\frac{J(J+1)-S(S+1)+K(K+1)}{2K(K+1)J(J+1)}.
\]

\ProblemHeading[3] The same problem for a ${}^1\Sigma$ term.

\SolutionHeading When $\Lambda=0$, the linear effect is absent, and second-order perturbation theory must be used. In carrying out the sum in the general formula (38.10), it is sufficient to retain only the terms corresponding to transitions between rotational components of the given electronic term (in the other terms, the energy differences in the denominators are large). We thus find
\[
\Delta E=d^2\mathcal E^2\left\{\frac{|\langle KM_K|n_z|K-1,M_K\rangle|^2}{E_K-E_{K-1}}+\frac{|\langle KM_K|n_z|K+1,M_K\rangle|^2}{E_K-E_{K+1}}\right\},
\]
where $E_K=BK(K+1)$. A straightforward calculation gives
\[
\Delta E=\frac{d^2\mathcal E^2}{B}\frac{K(K+1)-3M_K^2}{2K(K+1)(2K-1)(2K+3)}.
\]
